\documentclass[a4paper,12pt]{article}
\usepackage[utf8]{inputenc}
\usepackage[T1]{fontenc}
\usepackage[ngerman]{babel}
\usepackage{amsmath, amssymb}
\usepackage{geometry}
\usepackage{tikz}
\usetikzlibrary{shapes.geometric, arrows.meta, positioning, fit, backgrounds}
\usepackage{parskip}
\usepackage{titlesec}
\usepackage{booktabs}
\usepackage{enumitem}

\geometry{left=3cm, right=3cm, top=2.5cm, bottom=2.5cm}

\titleformat{\section}{\large\bfseries}{}{0em}{}[\titlerule]
\titleformat{\subsection}{\normalsize\bfseries}{}{0em}{}

\tikzset{
  box/.style={rectangle, draw, rounded corners=3pt, text width=4.5cm,
              align=center, minimum height=1.0cm, font=\small, inner sep=4pt},
  boxw/.style={box, text width=5.5cm},
  arrow/.style={-{Latex[length=2mm]}, thick},
}

\begin{document}

\begin{center}
  {\LARGE\bfseries Studierhinweise zu Lektion 6}\\[0.5em]
  {\large Mathematische Grundlagen (Modul 61111)}
\end{center}

\vspace{1em}

In dieser Lektion treten Sie in das Herzstück der Analysis ein. Kapitel~17 behandelt
\emph{höhere Ableitungen} differenzierbarer Funktionen. Im Mittelpunkt steht der
\textbf{Mittelwertsatz}, dessen einfacher Beweis weitreichende Folgerungen hat: Er
charakterisiert die Exponentialfunktion eindeutig, sichert die Eindeutigkeit der
Winkelfunktionen und liefert die \emph{Regel von de l'Hospital} zur Berechnung
schwieriger Grenzwerte. Der \textbf{Satz von Taylor} erlaubt es, beliebig oft
differenzierbare Funktionen durch Polynome zu approximieren und das Restglied
explizit abzuschätzen.

In Kapitel~18 werden \textbf{Reihen} als Folgen von Partialsummen eingeführt.
Da Grenzwerte von Reihen selten direkt berechenbar sind, stehen
\emph{Konvergenzkriterien} im Vordergrund: Majoranten-, Minoranten-,
Quotienten-, Wurzel- und Leibniz-Kriterium. Das Konzept der \emph{absoluten
Konvergenz} und der \emph{Potenzreihe} führt zu den sogenannten
\textbf{Summenfunktionen}, die stetig und sogar beliebig oft differenzierbar sind.

Kapitel~19 rundet die Lektion ab, indem Sinus, Cosinus, Tangens,
Arcusfunktionen, Exponentialfunktion, Logarithmus und allgemeine Potenzen als
\textbf{elementare Funktionen} sauber definiert und untersucht werden -- meist
über ihre Reihendarstellungen aus Kapitel~18.

Lesen Sie keinen Satz, ohne ihn verstanden zu haben. Zeichnen Sie Skizzen,
rechnen Sie Beispiele nach und lösen Sie die eingebetteten Aufgaben, bevor
Sie in die Lösungen schauen.

% -----------------------------------------------------------------------
\section*{Struktur der Lektion 6}

\begin{center}
\begin{tikzpicture}[node distance=0.55cm and 0.7cm]

  % Column headers
  \node[font=\bfseries\small] (h1) at (0,0)      {Kapitel 17: Höhere Ableitungen};
  \node[font=\bfseries\small] (h2) at (6.5,0)    {Kapitel 18: Reihen};
  \node[font=\bfseries\small] (h3) at (12.5,0)   {Kapitel 19: Elem.\ Funktionen};

  % Kap 17
  \node[box, below=0.8cm of h1] (mvs) {17.1 Mittelwertsatz\\(Satz von Rolle)};
  \node[box, below=0.5cm of mvs] (flg) {17.1 Folgerungen\\Monotonie, Extrema,\\de l'Hospital};
  \node[box, below=0.5cm of flg] (tay) {17.2 Taylorpolynome\\$P_{n,a}(x)$, Restglied};
  \node[box, below=0.5cm of tay] (sat) {17.2 Satz von Taylor\\Cauchy-/Lagrange-Form};

  % Kap 18
  \node[box, below=0.8cm of h2] (rei) {18.1 Reihen: Definition\\Partialsummen, Beispiele};
  \node[box, below=0.5cm of rei] (kon) {18.2 Konvergenzkriterien\\Majorante, Quotient,\\Wurzel, Leibniz};
  \node[box, below=0.5cm of kon] (abs) {18.3 Absolute Konvergenz\\Umordnungssatz};
  \node[box, below=0.5cm of abs] (pot) {18.4 Potenzreihen\\Konvergenzradius (Cauchy-Hadamard)};
  \node[box, below=0.5cm of pot] (sum) {18.5 Summenfunktionen\\stetig, differenzierbar};

  % Kap 19
  \node[box, below=0.8cm of h3] (sin) {19.1 sin und cos\\als Summenfunktionen};
  \node[box, below=0.5cm of sin] (kre) {19.1 Kreiszahl $\pi$,\\Periodizität, tan, cot};
  \node[box, below=0.5cm of kre] (zyc) {19.2 Zyklometrische\\Funktionen\\arcsin, arccos, arctan};
  \node[box, below=0.5cm of zyc] (exp) {19.3 exp, ln,\\allgemeine Potenz};

  % Kap 17 arrows
  \draw[arrow] (mvs) -- (flg);
  \draw[arrow] (flg) -- (tay);
  \draw[arrow] (tay) -- (sat);

  % Kap 18 arrows
  \draw[arrow] (rei) -- (kon);
  \draw[arrow] (kon) -- (abs);
  \draw[arrow] (abs) -- (pot);
  \draw[arrow] (pot) -- (sum);

  % Kap 19 arrows
  \draw[arrow] (sin) -- (kre);
  \draw[arrow] (kre) -- (zyc);
  \draw[arrow] (zyc) -- (exp);

  % Cross-chapter arrows
  \draw[arrow] (sat.east) -- ++(0.3,0) |- (rei.west);
  \draw[arrow] (sum.east) -- ++(0.3,0) |- (sin.west);

\end{tikzpicture}
\end{center}

\newpage

% -----------------------------------------------------------------------
\section*{Zielelement 17.1 -- Der Mittelwertsatz und seine Folgerungen}

\subsection*{Lerninhalte}

\begin{center}
\begin{tikzpicture}[node distance=0.6cm and 1.0cm]
  \node[box] (rol) {Satz von Rolle (17.1.2):\\$f(a)=f(b)$ $\Rightarrow$ $\exists\,x_0$: $f'(x_0)=0$};
  \node[box, below=0.5cm of rol] (mvs) {Mittelwertsatz (17.1.1):\\$f'(x_0)=\dfrac{f(b)-f(a)}{b-a}$};
  \node[box, below=0.5cm of mvs] (mon) {Monotonie aus Vorzeichen von $f'$\\(Korollar 17.1.8)};
  \node[box, below=0.5cm of mon] (ext) {Lokale Extrema:\\Korollare 17.1.9 und 17.1.10};
  \node[box, below=0.5cm of ext] (lhp) {Regel von de l'Hospital (17.1.14)\\$\lim \frac{f}{g} = \lim \frac{f'}{g'}$};
  \draw[arrow] (rol) -- (mvs);
  \draw[arrow] (mvs) -- (mon);
  \draw[arrow] (mon) -- (ext);
  \draw[arrow] (ext) -- (lhp);
\end{tikzpicture}
\end{center}

\subsection*{Lernziele}

Nach Durcharbeiten dieses Abschnitts sollten Sie:
\begin{itemize}
  \item den Satz von Rolle und den Mittelwertsatz kennen, geometrisch interpretieren
        und beweisen können,
  \item aus dem Vorzeichen von $f'$ auf das Monotonieverhalten von $f$ schließen
        können (Korollar~17.1.8),
  \item lokale Extrema mit den Kriterien 17.1.9 und 17.1.10 identifizieren können,
  \item die Regel von de l'Hospital anwenden können, inklusive des Spezialfalls 17.1.15,
        und wissen, wann sie \emph{nicht} anwendbar ist,
  \item die Folgerungen aus dem Mittelwertsatz verstehen: Eindeutigkeit der
        Exponentialfunktion (17.1.6) und der Winkelfunktionen (17.1.7).
\end{itemize}

\subsection*{Selbstkontrollelement 17.1}

Berechnen Sie $\displaystyle\lim_{x \to 0} \frac{1 - \cos(x)}{x^2}$ mit der Regel von de l'Hospital.
Begründen Sie, warum der Spezialfall (Korollar~17.1.15) beim ersten Schritt
\emph{nicht} direkt anwendbar ist.

\newpage

% -----------------------------------------------------------------------
\section*{Zielelement 17.2 -- Taylorpolynome und Satz von Taylor}

\subsection*{Lerninhalte}

\begin{center}
\begin{tikzpicture}[node distance=0.6cm and 1.0cm]
  \node[box] (koef) {Koeffizienten $a_k = \dfrac{f^{(k)}(a)}{k!}$};
  \node[box, below=0.5cm of koef] (def) {Taylorpolynom $P_{n,a}(x)$ (Def.~17.2.1)};
  \node[box, below=0.5cm of def] (rest) {Restglied $R_{n,a}(x) = f(x) - P_{n,a}(x)$};
  \node[box, below=0.5cm of rest] (form) {Cauchy-Form und Lagrange-Form\\des Restglieds (17.2.10)};
  \node[box, below=0.5cm of form] (absch) {Restgliedabschätzung (17.2.12):\\$|R_{n,a}(x)| \le \dfrac{K}{(n+1)!}|x-a|^{n+1}$};
  \node[box, below=0.5cm of absch] (bsp) {Beispiele: exp, sin, cos, ln};
  \draw[arrow] (koef) -- (def);
  \draw[arrow] (def) -- (rest);
  \draw[arrow] (rest) -- (form);
  \draw[arrow] (form) -- (absch);
  \draw[arrow] (absch) -- (bsp);
\end{tikzpicture}
\end{center}

\subsection*{Lernziele}

Nach Durcharbeiten dieses Abschnitts sollten Sie:
\begin{itemize}
  \item Taylorpolynome einer Funktion $f$ in einem Entwicklungspunkt $a$ berechnen können,
  \item die beiden Formen des Restglieds (Cauchy und Lagrange) kennen und im
        Beweis des Satzes von Taylor einordnen können,
  \item Restglieder mithilfe von Korollar~17.2.12 abschätzen können,
  \item verstehen, wann die Taylorentwicklung gegen $f(x)$ konvergiert
        (Proposition~17.2.14),
  \item die Taylorpolynome der elementaren Funktionen $\exp$, $\sin$, $\cos$ und
        $\ln(1+x)$ im Entwicklungspunkt 0 angeben können.
\end{itemize}

\subsection*{Selbstkontrollelement 17.2}

Bestimmen Sie das 4.\ Taylorpolynom $P_{4,0}$ von $f(x) = \cos(x)$ im
Entwicklungspunkt $a = 0$ und geben Sie eine Abschätzung für das Restglied
$|R_{4,0}(x)|$ auf dem Intervall $[-1, 1]$.

\newpage

% -----------------------------------------------------------------------
\section*{Zielelement 18.1 -- Reihen: Grundlagen}

\subsection*{Lerninhalte}

\begin{center}
\begin{tikzpicture}[node distance=0.6cm and 1.0cm]
  \node[box] (def) {Reihe $\sum_{n=1}^{\infty} a_n$:\\Folge der Partialsummen $(s_n)$};
  \node[box, below=0.5cm of def] (kon) {Konvergenz: $\lim s_n = s$};
  \node[box, below=0.5cm of kon] (bsp) {Beispiele: geom.\ Reihe,\\harmonische Reihe,\\alternierende harm.\ Reihe};
  \node[box, below=0.5cm of bsp] (tay) {Taylorreihe $\sum_{k=0}^{\infty}\frac{f^{(k)}(a)}{k!}(x-a)^k$\\und ihre Konvergenz};
  \node[box, below=0.5cm of tay] (rec) {Rechenregeln (18.1.20):\\Summe, skalare Multiplikation};
  \draw[arrow] (def) -- (kon);
  \draw[arrow] (kon) -- (bsp);
  \draw[arrow] (bsp) -- (tay);
  \draw[arrow] (tay) -- (rec);
\end{tikzpicture}
\end{center}

\subsection*{Lernziele}

Nach Durcharbeiten dieses Abschnitts sollten Sie:
\begin{itemize}
  \item den Begriff der Reihe als Folge von Partialsummen erklären können,
  \item wissen, dass aus der Konvergenz einer Reihe folgt, dass $(a_n)$ eine
        Nullfolge ist (Proposition~18.1.16), und das Umgekehrte widerlegen können,
  \item die geometrische Reihe kennen und für $|q| < 1$ ihren Grenzwert angeben,
  \item den Zusammenhang zwischen Taylorreihe und Restglied kennen
        (Proposition~18.1.11),
  \item Exponential-, Logarithmus- und alternierende harmonische Reihe als
        Taylorreihen einordnen.
\end{itemize}

\subsection*{Selbstkontrollelement 18.1}

Zeigen Sie, dass $\displaystyle\sum_{n=1}^{\infty} \frac{1}{n(n+1)} = 1$ gilt, indem Sie
die Partialsummen explizit berechnen (Hinweis: Partialbruchzerlegung
$\frac{1}{n(n+1)} = \frac{1}{n} - \frac{1}{n+1}$).

\newpage

% -----------------------------------------------------------------------
\section*{Zielelement 18.2 -- Konvergenzkriterien für Reihen}

\subsection*{Lerninhalte}

\begin{center}
\begin{tikzpicture}[node distance=0.6cm and 1.0cm]
  \node[box] (maj) {Majorantenkriterium (18.2.1):\\$|a_n| \le b_n$, $\sum b_n$ konvergent};
  \node[box, right=0.8cm of maj] (min) {Minorantenkriterium (18.2.6):\\$a_n \ge b_n$, $\sum b_n$ divergent};
  \node[box, below=0.5cm of maj] (quo) {Quotientenkriterium (18.2.8):\\$\frac{|a_{n+1}|}{|a_n|} \le q < 1$};
  \node[box, below=0.5cm of min] (wur) {Wurzelkriterium (18.2.11):\\$\sqrt[n]{|a_n|} \le q < 1$};
  \node[box, below=0.5cm of quo] (gre) {Korollar 18.2.14:\\Grenzwert $< 1$ oder $> 1$};
  \node[box, below=0.5cm of wur] (lei) {Leibniz-Kriterium (18.2.24):\\$(b_n)$ monoton fallende Nullfolge};
  \draw[arrow] (maj) -- (quo);
  \draw[arrow] (min) -- (wur);
  \draw[arrow] (quo) -- (gre);
  \draw[arrow] (wur) -- (lei);
\end{tikzpicture}
\end{center}

\subsection*{Lernziele}

Nach Durcharbeiten dieses Abschnitts sollten Sie:
\begin{itemize}
  \item die fünf Konvergenzkriterien (Majorante, Minorante, Quotient, Wurzel,
        Leibniz) kennen, deren Voraussetzungen prüfen und anwenden können,
  \item wissen, warum $\frac{|a_{n+1}|}{|a_n|} < 1$ für alle $n$ noch keine
        Konvergenz garantiert (Gegenbeispiel: harmonische Reihe),
  \item das handliche Korollar~18.2.14 (Grenzwert von Quotienten- oder
        Wurzelfolge) kennen und einsetzen können,
  \item Cosinusreihe, Sinusreihe und Binomialreihe mit dem Quotientenkriterium auf
        Konvergenz untersuchen können.
\end{itemize}

\subsection*{Selbstkontrollelement 18.2}

Untersuchen Sie die Reihe $\displaystyle\sum_{n=1}^{\infty} \frac{n}{2^n}$ mit dem
Quotientenkriterium (Korollar~18.2.14) auf Konvergenz.

\newpage

% -----------------------------------------------------------------------
\section*{Zielelement 18.3--18.4 -- Absolute Konvergenz und Potenzreihen}

\subsection*{Lerninhalte}

\begin{center}
\begin{tikzpicture}[node distance=0.6cm and 1.0cm]
  \node[box] (abs) {Absolute Konvergenz (18.3.1):\\$\sum |a_n|$ konvergiert};
  \node[box, below=0.5cm of abs] (imp) {Absolut konvergent $\Rightarrow$ konvergent\\(Proposition 18.3.4)};
  \node[box, below=0.5cm of imp] (umo) {Umordnung abs.\ konvergenter\\Reihen (18.3.6)};
  \node[box, below=0.5cm of umo] (pot) {Potenzreihe $\sum a_n x^n$ (18.4.1)\\Konvergenzradius $R$};
  \node[box, below=0.5cm of pot] (ch)  {Cauchy-Hadamard (18.4.13):\\$R = \dfrac{1}{\limsup \sqrt[n]{|a_n|}}$};
  \draw[arrow] (abs) -- (imp);
  \draw[arrow] (imp) -- (umo);
  \draw[arrow] (umo) -- (pot);
  \draw[arrow] (pot) -- (ch);
\end{tikzpicture}
\end{center}

\subsection*{Lernziele}

Nach Durcharbeiten dieses Abschnitts sollten Sie:
\begin{itemize}
  \item den Begriff der absoluten Konvergenz erklären und von der bloßen
        Konvergenz abgrenzen können,
  \item den Umordnungssatz (18.3.6) und den Satz von Riemann (18.3.7)
        kennen und einordnen können,
  \item den Begriff des Konvergenzradius verstehen: für $|x| < R$ absolute
        Konvergenz, für $|x| > R$ Divergenz (Proposition~18.4.7),
  \item den Satz von Cauchy-Hadamard anwenden können,
  \item die Konvergenzradien der wichtigsten Potenzreihen (Tabelle~18.5.1) kennen.
\end{itemize}

\subsection*{Selbstkontrollelement 18.3--18.4}

Bestimmen Sie den Konvergenzradius der Potenzreihe
$\displaystyle\sum_{n=0}^{\infty} \frac{x^n}{n!}$ mit dem Satz von Cauchy-Hadamard.

\newpage

% -----------------------------------------------------------------------
\section*{Zielelement 18.5 -- Summenfunktionen}

\subsection*{Lerninhalte}

\begin{center}
\begin{tikzpicture}[node distance=0.6cm and 1.0cm]
  \node[box] (def) {Summenfunktion $f(x) = \sum a_n x^n$\\auf Konvergenzintervall $K$};
  \node[box, below=0.5cm of def] (ste) {Stetigkeit (18.5.4)};
  \node[box, below=0.5cm of ste] (dif) {Differenzierbarkeit (18.5.5):\\$f'(x) = \sum n\,a_n x^{n-1}$};
  \node[box, below=0.5cm of dif] (inf) {Unendlich oft diffbar.\ (18.5.6)};
  \node[box, below=0.5cm of inf] (stm) {Stammfunktion (18.5.11):\\$F(x) = \sum \frac{a_n}{n+1} x^{n+1}$};
  \draw[arrow] (def) -- (ste);
  \draw[arrow] (ste) -- (dif);
  \draw[arrow] (dif) -- (inf);
  \draw[arrow] (inf) -- (stm);
\end{tikzpicture}
\end{center}

\subsection*{Lernziele}

Nach Durcharbeiten dieses Abschnitts sollten Sie:
\begin{itemize}
  \item wissen, dass Summenfunktionen stetig und unendlich oft differenzierbar sind,
  \item die gliedweise Differentiation einer Potenzreihe anwenden können,
  \item verstehen, dass jede Potenzreihe die Taylorreihe ihrer Summenfunktion ist
        (Korollar~18.5.8),
  \item eine Stammfunktion einer Summenfunktion durch gliedweise Integration
        bestimmen können (18.5.11),
  \item die Logarithmusreihe und die Binomialreihe mithilfe dieser Methoden auf
        ganz $(-1, 1)$ ausdehnen (Propositions 18.5.12 und 18.5.14).
\end{itemize}

\subsection*{Selbstkontrollelement 18.5}

Zeigen Sie durch gliedweise Differentiation der Exponentialreihe
$\exp(x) = \sum_{n=0}^{\infty} \frac{x^n}{n!}$, dass $\exp' = \exp$ gilt.

\newpage

% -----------------------------------------------------------------------
\section*{Zielelement 19.1--19.2 -- Trigonometrische und zyklometrische Funktionen}

\subsection*{Lerninhalte}

\begin{center}
\begin{tikzpicture}[node distance=0.6cm and 1.0cm]
  \node[box] (def) {sin und cos als Summenfunktionen\\(Sinusreihe, Cosinusreihe)};
  \node[box, below=0.5cm of def] (eig) {Eigenschaften (SinCos-1) bis (SinCos-5):\\stetig, gerade/ungerade, Additionstheoreme};
  \node[box, below=0.5cm of eig] (py)  {Trig.\ Pythagoras: $\sin^2+\cos^2=1$\\$|\sin|, |\cos| \le 1$};
  \node[box, below=0.5cm of py]  (pi)  {Kreiszahl $\pi$ (Def.~19.1.6),\\spezielle Werte, Periodizität $2\pi$};
  \node[box, below=0.5cm of pi]  (tan) {tan, cot:\\Definitionsbereiche, Ableitungen,\\Monotonie, Periode $\pi$};
  \node[box, below=0.5cm of tan] (zyc) {Zyklometrische Funktionen:\\arcsin, arccos, arctan, arccot,\\ihre Ableitungen};
  \draw[arrow] (def) -- (eig);
  \draw[arrow] (eig) -- (py);
  \draw[arrow] (py)  -- (pi);
  \draw[arrow] (pi)  -- (tan);
  \draw[arrow] (tan) -- (zyc);
\end{tikzpicture}
\end{center}

\subsection*{Lernziele}

Nach Durcharbeiten dieses Abschnitts sollten Sie:
\begin{itemize}
  \item Sinus und Cosinus über ihre Reihendarstellungen definieren und die
        Eigenschaften aus Satz~19.1.11 vollständig kennen,
  \item die Kreiszahl $\pi$ als Nullstelle des Cosinus definieren und die
        speziellen Funktionswerte auswendig kennen,
  \item die Additionstheoreme und den trigonometrischen Pythagoras anwenden,
  \item Tangens und Cotangens als Quotienten erklären, ihre Ableitungen herleiten
        und die Periodizität begründen können,
  \item die zyklometrischen Funktionen als Umkehrfunktionen auf geeigneten
        Intervallen erklären und ihre Ableitungen kennen (arcsin$'(x) = \frac{1}{\sqrt{1-x^2}}$,
        arctan$'(x) = \frac{1}{1+x^2}$).
\end{itemize}

\subsection*{Selbstkontrollelement 19.1--19.2}

Beweisen Sie $\sin(2x) = 2\sin(x)\cos(x)$ mithilfe der Additionstheoreme und
bestimmen Sie anschließend $\arctan(1)$ und $\arctan(-1)$.

\newpage

% -----------------------------------------------------------------------
\section*{Zielelement 19.3 -- Exponentialfunktion, Logarithmus und allgemeine Potenz}

\subsection*{Lerninhalte}

\begin{center}
\begin{tikzpicture}[node distance=0.6cm and 1.0cm]
  \node[box] (exp) {Exponentialfunktion $\exp$:\\Reihendarstellung, exp$' = $ exp,\\$\exp(x+y) = \exp(x)\cdot\exp(y)$};
  \node[box, below=0.5cm of exp] (ln)  {Natürlicher Logarithmus $\ln$:\\Umkehrfunktion von exp,\\$\ln'(x) = 1/x$};
  \node[box, below=0.5cm of ln]  (pot) {Allgemeine Potenz $x^\alpha$:\\$f'(x) = \alpha x^{\alpha-1}$};
  \node[box, below=0.5cm of pot] (zus) {Zusammenhang:\\$x^a = \exp(a\ln x)$,\\$\exp_a(x) = \exp(x\ln a)$,\\$\log_a(x) = \frac{\ln x}{\ln a}$};
  \draw[arrow] (exp) -- (ln);
  \draw[arrow] (ln)  -- (pot);
  \draw[arrow] (pot) -- (zus);
\end{tikzpicture}
\end{center}

\subsection*{Lernziele}

Nach Durcharbeiten dieses Abschnitts sollten Sie:
\begin{itemize}
  \item die Exponentialfunktion über ihre Reihendarstellung definieren und die
        Eigenschaften aus Satz~19.3.1 kennen,
  \item den natürlichen Logarithmus als Umkehrfunktion von $\exp$ erklären,
        seine Reihendarstellung und Ableitungsregel $\ln'(x) = 1/x$ kennen,
  \item die allgemeine Potenz $x^\alpha$ für $\alpha \in \mathbb{R}$ verstehen und
        die Ableitungsregel anwenden können,
  \item den Zusammenhang zwischen allgemeiner Potenz, Exponentialfunktion und
        Logarithmus zur Basis $a$ aus Proposition~19.3.4 kennen.
\end{itemize}

\subsection*{Selbstkontrollelement 19.3}

Berechnen Sie $\dfrac{d}{dx} x^{\pi}$ und $\dfrac{d}{dx} \pi^x$ für $x > 0$,
wobei $\pi$ die Kreiszahl ist. Erklären Sie den Unterschied zwischen
$x^\pi$ (allgemeine Potenz) und $\pi^x$ (Exponentialfunktion zur Basis $\pi$).

\end{document}
